\documentclass{article}
\usepackage[spanish]{babel}
\usepackage[utf8]{inputenc}
\usepackage{enumitem}
\usepackage{amsmath}
\usepackage{amssymb}
\usepackage{amsthm}
\usepackage{geometry}
\usepackage{setspace}
\usepackage{fancyhdr}
\geometry{letterpaper, margin=.7in}

% Edita rápido el ejercicio y la tarea
\newcommand{\ntarea}{1}

\begin{document}

\pagestyle{fancy}
\lhead{Probabilidad}
\rhead{Tarea \ntarea}

\large{Jesús Ezequiel Ramos Moreno.}
\begin{enumerate}
    \item[2.] Demuestre que si \(\{A_n\}\) es una colección de eventos en un mismo espacio de probabilidad todos con probabilidad 1, entonces \(\mathbb{P}(\bigcap_{n=1}^{\infty} A_n) = 1\).
    \begin{proof}
        Consideremos los complementos. Para cada \(n\), como \(P(A_n) = 1\), entonces
        \[
        P(A_n^c) = 1 - P(A_n) = 0.
        \]
        Por las leyes de De Morgan,
        \[
        {\left(\bigcap_{n = 1}^\infty A_n\right)}^c = \bigcup_{n = 1}^\infty A_n ^c.
        \]
        Por la propiedad de la subaditividad numerable de la medida \(P\), tenemos que
        \[
        P\left(\bigcup_{n = 1} ^\infty A_n ^c\right) \leq \sum_{n = 1}^{\infty}P(A_n^c) = \sum_{n=1}^{\infty}0 = 0.
        \]
        Como la propiedad es no negativa, \(P\left(\bigcup_{n = 1} ^\infty A_n ^c\right) = 0\). Entonces 
        \[
        P\left(\bigcap_{i = 1}^\infty A_n\right) = 1 - P\left(\bigcup_{n = 1} ^\infty A_n ^c\right) = 1 - 0 = 1.
        \]
    \end{proof}
    \item[4.] Sean \(F(x)\) y \(G(x)\) dos funciones de distribución. Determine si las siguientes funciones son de distribución.
    \begin{enumerate}[label={(\roman*)}]
        \item \(aF(x) + (1 - a)G(x)\), con \(0 \leq a \leq 1\)
        \item \(F(x) + G(x)\)
        \item \(F(x)G(x)\)
        \item \(\frac{2G(x)}{1 + F(x)}\)
    \end{enumerate}
    \textbf{Solución.} Verificaremos las propiedades de una función de distribución para cada inciso.
    \begin{enumerate}[label={(\roman*)}]
        \item Notemos que cuando \(a\) es igual a 0 o 1, \(aF(x) + (1 - a)G(x)\) es igual a \(G(x)\) o \(F(x)\), respectivamente, por lo que es una función de distribución.
        
        Consideremos \(a \in (0, 1)\) y verifiquemos las propiedades de la función.
        \begin{itemize}
            \item Calculemos los límites al infinito: 
            \[
            \lim_{x \to -\infty} aF(x) + (1 - a)G(x) \quad \text{y} \lim_{x \to +\infty} aF(x) + (1 - a)G(x).
            \]
            Ya que \(F\) y \(G\) son funciones de distribución, su límite cuando \(x\) tiende a infinito negativo es 0, por lo que al evaluar el límite tenemos:
            \[
            \lim_{x \to -\infty} aF(x) + (1 - a)G(x) = a(0) + (1 - a)(0) = 0.
            \]
            Análogamente, cuando \(x\) tiende a infinito positivo, el límite de \(F\) y \(G\) es 1. Por lo tanto,
            \[
            \lim_{x \to +\infty} aF(x) + (1 - a)G(x) = a(1) + (1 - a)(1) = a + 1 - a = 1.
            \]
            \item Dado que \(a \geq 0\), entonces la monotonicidad del primer término \(aF(x)\) es igual a la de \(F(x)\), que es creciente por hipótesis. De forma similar, como \(a \leq 1\), \((1 - a) \geq 0\), por lo que el término \((1 - a)G(x)\) es creciente también. Dado que la suma de funciones crecientes es creciente, se sigue que la función \(aF(x) + (1 - a)G(x)\) es creciente.
            \item Ya que hemos probado que la función es creciente y su valor en sus límites es 0 y 1, se sigue que 
            \[
            0 \leq aF(x) + (1 - a)G(x) \leq 1.
            \]
        \end{itemize}
        Ya que hemos verificado las propiedades, entonces la función \(aF(x) + (1 - a)G(x)\) es una función de distribución.
        \item Veamos que \(F(x) + G(x)\) no es una función de distribución. Calculemos el límite cuando \(x\) tiende a infinito positivo. Dado que \(F\) y \(G\) son funciones de distribución, ambas convergen a 1 cuando \(x\) tiende a infinito. Evaluando el límite, obtenemos 
    \[
    \lim_{x \to +\infty} [F(x) + G(x)] = 1 + 1 = 2,
    \]
    lo cuál viola las propiedades de una función de distribución. Por lo tanto, \(F(x) + G(x)\) no es una función de distribución.
    \item \begin{itemize}
        \item Calculemos los límites al infinito:
        \[
        \lim_{x \to -\infty}F(x)G(x) \quad \text{y} \quad \lim_{x \to +\infty} F(x)G(x).
        \]
    \end{itemize}
    \item a
\end{enumerate}
\end{enumerate}

\end{document}